\chapter{Background Study}

This chapter establishes the theoretical and empirical foundation of FEDMED-X.
First, it reviews the major research streams relevant to federated medical
image classification: non-IID robust optimization, privacy-preserving
learning, communication-efficient model updates, and edge-oriented deployment.
Second, it analyzes the practical problem space and identifies the systems
gap that motivates the proposed framework. The goal of this chapter is not
only to summarize prior work, but to show why isolated solutions are
insufficient for real hospital constraints and why a joint evaluation
pipeline is required.

\section{Literature Review}
Federated Learning (FL) has become a central paradigm for collaborative
healthcare AI because it enables multi-institutional training without direct
exchange of raw patient images~\cite{mcmahan2017communication,
guan2024federated}. In medical imaging, this property is especially important
due to strict legal and ethical constraints on data movement. Over the last
several years, the literature has expanded from basic aggregation methods
toward specialized solutions for statistical heterogeneity, formal privacy,
communication bottlenecks, and deployment realism.

\subsection{Federated Learning for Medical Imaging}
Guan and Liu provide a comprehensive survey of federated medical image
analysis and organize the field around system architecture, privacy
mechanisms, and communication challenges~\cite{guan2024federated}. Their
review makes clear that medical FL is not a single-algorithm problem; it is a
multi-constraint systems problem. Ghosh \emph{et al.} further emphasize that
despite rapid publication growth, clinical translation remains limited and
deployment-oriented benchmarks are still weak~\cite{ghosh2026advancements}.
Zhou \emph{et al.} reinforce this point by arguing that fair evaluation of
medical FL should include not only accuracy, but also communication and
computational efficiency~\cite{zhou2025federated}. Hindawi \emph{et al.}
provide a multi-center feasibility study for lesion segmentation and confirm
the operational complexity of real hospital
federations~\cite{hindawi2025federated}. These surveys collectively indicate
a shift in research expectations. Modern medical FL research is judged by
whether it remains robust under non-IID partitions, whether it supports
privacy constraints, and whether it can operate under practical resource
limits.

\subsection{Non-IID Heterogeneity and FedProx}
A major obstacle in multi-hospital learning is non-IID data. Different
institutions observe different class distributions, acquisition protocols,
and patient demographics. Under such heterogeneity, standard Federated
Averaging (FedAvg) often suffers from client drift and unstable convergence.
Li \emph{et al.} introduced FedProx to address heterogeneous federated
optimization by adding a proximal term that penalizes local deviation from
the global model~\cite{li2020fedprox}. Yuan and Li later derived stronger
convergence guarantees for FedProx under local dissimilarity
conditions~\cite{yuan2022convergence}. Alekseenko \emph{et al.} proposed a
distance-aware alternative that downweights distant clients in medical
segmentation~\cite{alekseenko2024distance}, and Maleki and Khotanlou combined
Vision Transformers with SCAFFOLD to handle strong heterogeneity in medical
imaging~\cite{maleki2025handling}.

Riyadi \emph{et al.} provide direct empirical support for FedProx in medical
imaging benchmarks. In a comparative study on MedMNIST tasks under Dirichlet
label skew, FedProx demonstrated stronger robustness and stability than
FedAvg and other selected baselines~\cite{riyadi2026comparative}. However,
Riyadi \emph{et al.} do not jointly evaluate formal differential privacy,
aggressive update compression, and edge deployment constraints. Therefore,
their results establish FedProx as a strong non-IID baseline, but not as a
complete deployment solution.

\subsection{Privacy Risks and Differential Privacy}
A common misconception in federated healthcare systems is that data locality
automatically guarantees privacy. In reality, shared gradients and model
updates can still leak sensitive information through reconstruction and
membership-inference style attacks. Koutsoubis \emph{et al.} highlight this
residual risk~\cite{koutsoubis2024future}, and Wei \emph{et al.} recently
demonstrated concrete gradient-inversion attacks against sparsified medical
FL~\cite{wei2025medical}. Abadi \emph{et al.} established Differentially
Private Stochastic Gradient Descent (DP-SGD) as a foundational method for
training deep networks with formal privacy guarantees by clipping gradients
and injecting calibrated noise~\cite{abadi2016deep}, with tighter accounting
provided by R\'enyi Differential Privacy~\cite{mironov2017renyi}.

In the medical FL literature, privacy-preserving methods vary widely in
strength and practicality. Ziller \emph{et al.} applied DP to multi-site
medical image segmentation and quantified the utility
cost~\cite{ziller2021differentially}. Hou \emph{et al.} studied the joint
tuning of local steps and communication rounds under a fixed privacy
budget~\cite{hou2023improving}, while Yu \emph{et al.} proposed adaptive
gradient-component privacy allocation to reduce noise
waste~\cite{yu2024adaptive}. Jiang \emph{et al.} recently introduced a
shadow-model defense against targeted gradient-inversion attacks on chest
X-ray data~\cite{jiang2025shadow}. Yesmin proposed MedHE, combining
communication efficiency with privacy-preserving mechanisms including DP and
homomorphic encryption~\cite{yesmin2025medhe}. Pachlasiya and Patel reported
private multi-hospital medical imaging experiments with explicit privacy
budgeting~\cite{pachlasiya2026fedhealth}. These studies confirm that privacy
can be integrated into medical FL, but they also reveal a recurring
limitation: privacy-preserving pipelines are rarely evaluated together with
edge-hardware constraints and complete bandwidth--latency trade-offs.

For FEDMED-X, DP-SGD is selected because it provides a mathematically
interpretable privacy budget $(\epsilon,\delta)$ and can be implemented in
practical deep learning pipelines through frameworks such as Opacus. This
choice prioritizes measurable privacy--utility analysis over heavier
cryptographic alternatives that may be difficult to deploy on low-resource
clinical edge devices.

\subsection{Communication Efficiency and Compression}
Communication overhead is one of the most immediate barriers to real hospital
FL. Model updates can be large, rounds are repeated many times, and
institutional networks may be restricted. Schoenpflug \emph{et al.} document
how network constraints and operational complexity affect real multi-center
FL deployments~\cite{schoenpflug2025navigating}. In response, multiple
studies investigate sparsification, quantization, and other compression
strategies. Zhou \emph{et al.} combined a proximal term with adaptive model
transmission for diabetic retinopathy classification~\cite{zhou2022communication}.
Asad and Yuan proposed FeSEC, a joint secure and communication-efficient
framework that reduces payloads below 10\% of a FedAvg
comparator~\cite{asad2024fesec}. Habib \emph{et al.} explicitly transmit only
TopK gradients for medical FL~\cite{habib2025cflsparsemed}, while Sun and Sun
reported low-rank federated distillation with strong communication
savings~\cite{sun2025medfd}. Yan \emph{et al.} took an alternative route
using condensed knowledge distillation for non-IID medical
FL~\cite{yan2024fedvck}. These works demonstrate that bandwidth can be
reduced significantly, but compression is not free. Aggressive sparsification
can discard informative update components, and when combined with DP noise,
the interaction may further degrade utility. This is exactly why FEDMED-X
treats compression as an experimental factor rather than an assumed
improvement.

\subsection{Edge Deployment and Clinical Translation}
Even when FL algorithms perform well in simulation, clinical value depends on
deployability. Bujotzek \emph{et al.} show that real radiology FL initiatives
face infrastructure and process barriers beyond model
accuracy~\cite{bujotzek2024real}. Vincent and Surantha investigate federated
deep learning and real-time inference on edge computing devices for skin
cancer classification, providing evidence that edge-oriented medical imaging
workflows are feasible~\cite{vincent2026federated}. Song \emph{et al.}
developed a cloud--edge split federated platform on the Jetson Nano for
breast-cancer pathology~\cite{song2025edgecloud}. Shabarinath \emph{et al.}
benchmarked ResNet-18 transfer learning on the Jetson Nano for pneumonia,
tuberculosis, and lung-cancer classification, reporting per-image latency in
the millisecond range~\cite{shabarinath2025efficacy}. Eshwarappa \emph{et
al.} showed that Edge-Avg federated learning reduces training latency by up
to 47\% in a multi-institutional skin cancer
setting~\cite{eshwarappa2025communication}. Memon \emph{et al.} evaluated
Jetson Xavier NX, Raspberry Pi, and Android platforms for brain MRI
classification and segmentation, quantifying accuracy--latency
trade-offs~\cite{memon2024edge}.

Nevertheless, edge papers often focus on inference speed or device
feasibility and do not always integrate formal DP accounting, non-IID robust
optimization, and compression analysis in one pipeline. This separation of
concerns is the key literature gap. Ghosh \emph{et al.} explicitly note the
shortage of deployment-oriented evidence in medical FL
research~\cite{ghosh2026advancements}. Zhou \emph{et al.} similarly call for
broader evaluation criteria that include system cost
dimensions~\cite{zhou2025federated}. FEDMED-X is designed in direct response
to this gap.

\subsection{Lightweight Medical Models and Attention Mechanisms}
Deployment on resource-constrained hardware requires compact models. Large
backbones may achieve strong accuracy in cloud settings but are less suitable
for low-power clinical edge devices. For this reason, FEDMED-X uses a
lightweight custom MedNet architecture. To improve representational focus on
diagnostically relevant regions, the model incorporates the Convolutional
Block Attention Module (CBAM), which combines channel and spatial
attention~\cite{woo2018cbam}. In addition, class imbalance is common in
medical datasets; Focal Loss is therefore used to emphasize difficult
samples during training~\cite{lin2017focal}. The dataset foundation of this
project is MedMNIST v2, a standardized biomedical image classification
benchmark that supports reproducible experimentation~\cite{yang2023medmnist}.
Although MedMNIST images are smaller and more controlled than full-resolution
clinical scans, they provide a valid and widely used testbed for
methodological comparison and systems ablation under federated partitions.

\subsection{Synthesis of the Literature}
Table~\ref{tab:litgap} summarizes the conceptual gap using representative
prior directions. The pattern is consistent: existing research provides
strong partial solutions, but joint evaluation under non-IID, privacy,
compression, and edge constraints remains limited.

\begin{table}[!htbp]
\centering
\caption{Conceptual comparison of prior research directions and FEDMED-X focus.}
\label{tab:litgap}
\small
\begin{tabularx}{\textwidth}{@{}l X X c@{}}
\toprule
\textbf{Research Direction} & \textbf{Prior Contribution} & \textbf{Typical Limitation} & \textbf{Addressed in FEDMED-X} \\
\midrule
Medical FL surveys        & Map privacy, architecture, and communication issues       & Often no complete system implementation & Yes \\
FedProx / non-IID FL      & Improve robustness under heterogeneous clients            & Limited joint privacy + compression + edge evaluation & Yes \\
DP medical FL             & Formal or practical privacy protection                    & Weak edge deployment analysis            & Yes \\
Compressed FL             & Reduce communication cost                                 & Privacy interaction not always quantified & Yes \\
Edge medical AI           & Show on-device feasibility                                & Often inference-only or incomplete FL stack & Partially \\
\textbf{FEDMED-X}         & Unified FedProx + DP-SGD + TopK/INT8 + MedNet-CBAM        & Controlled benchmark setting             & \textbf{Core focus} \\
\bottomrule
\end{tabularx}
\end{table}

\noindent\emph{Table Explanation:} Table~\ref{tab:litgap} contrasts five
prior research directions against the unified scope of FEDMED-X. Each
individual direction addresses one aspect of the medical FL problem but
leaves the other three uncovered, which motivates the systems-integration
contribution of this work.
\FloatBarrier

\section{Problem Analysis}
The literature makes clear that the core problem is multi-dimensional.
FEDMED-X therefore analyzes the problem at four interdependent layers.

\subsection{Statistical Problem: Non-IID Hospital Data}
In multi-hospital learning, local datasets differ in size and class balance.
Under Dirichlet partitioning, some clients may contain dominant
disease-positive samples while others contain mostly normal cases. This
imbalance causes local optima to diverge. Standard FedAvg aggregates these
divergent updates without sufficiently constraining local drift, which can
reduce global stability. FedProx addresses this statistical problem by
penalizing deviation from the global model through a proximal
term~\cite{li2020fedprox,riyadi2026comparative,yuan2022convergence}.

\subsection{Security Problem: Residual Gradient Leakage}
Even if raw images never leave the hospital, shared updates may leak private
information~\cite{koutsoubis2024future,wei2025medical}. From a
systems-security perspective, data locality is necessary but not sufficient.
Formal privacy mechanisms such as DP-SGD are required when regulatory
confidence and adversarial robustness matter~\cite{abadi2016deep,
ziller2021differentially}. The practical challenge is that stronger privacy
noise usually reduces diagnostic utility. Hence privacy is a tunable
trade-off that must be measured, not a binary feature.

\subsection{Systems Problem: Communication and Network Constraints}
Hospital networks are not experimental testbeds. Firewalls, limited
bandwidth, and operational controls can delay or block heavy update
exchange~\cite{bujotzek2024real,schoenpflug2025navigating}. Compression
techniques such as TopK sparsification and low-precision quantization reduce
payload size, but may also remove useful signal~\cite{habib2025cflsparsemed,
asad2024fesec}. When compression is combined with DP noise, discarded
coefficients may include information needed for recovery of model quality.
This creates a joint optimization problem rather than an independent add-on
feature.

\subsection{Deployment Problem: Edge Resource Limits}
A clinically useful pipeline must eventually run in constrained environments.
Rural and semi-urban clinics cannot always depend on continuous cloud
connectivity or high-end GPUs. This creates demand for lightweight models,
efficient inference, and deployment-aware architecture
choices~\cite{vincent2026federated,shabarinath2025efficacy,memon2024edge}.
FEDMED-X therefore treats edge suitability as a design requirement, not a
postscript.

\subsection{Integrated Problem Statement}
Combining the four layers yields the integrated problem addressed by this
project:

\begin{quote}\itshape
How can a federated medical image classification system preserve non-IID
robustness, enforce formal privacy, reduce communication cost, and remain
eligible for edge deployment, while quantifying the resulting accuracy
trade-offs under a controlled and reproducible experimental protocol?
\end{quote}

FEDMED-X answers this question through systems integration and factorial
evaluation rather than through a claim of a single new optimizer. The
expected contribution is a practical evidence base that shows what
performance remains when all major deployment constraints are applied
together.

\subsection{Implications for Methodology}
This problem analysis directly determines the methodological choices in
Chapter~3:
\begin{enumerate}[leftmargin=*]
  \item use FedProx for non-IID stability;
  \item use DP-SGD for formal privacy accounting;
  \item use TopK and INT8 for communication reduction;
  \item use MedNet-CBAM for lightweight representation learning;
  \item evaluate all components through controlled ablation on MedMNIST
        federated partitions.
\end{enumerate}

\chaptersummary{This chapter reviewed the background necessary to understand
FEDMED-X. Prior research provides strong foundations in medical federated
learning, FedProx-based non-IID optimization, differential privacy,
communication compression, and edge-oriented inference. However, these
components are usually studied in separation. Clinical deployment studies and
surveys confirm that real hospital environments impose joint constraints
that isolated benchmarks do not capture. The resulting research gap is a
lack of unified, reproducible systems evaluation under non-IID, privacy,
bandwidth, and edge-related requirements. Chapter~3 converts this gap into a
concrete methodology for the FEDMED-X framework.}